% =====================================================================
% INTRODUCCION -- SECCION A CARGO DE OTRO INTEGRANTE
% Responsable: no corresponde al frente etico-legal.
% Fabian: no editar el contenido; mantener el \section y los \label para
% que las referencias cruzadas del resto del informe no se rompan.
%
% ESTADO (ver PENDIENTES.md): la decision de arquitectura (B-01/D-11) sigue
% abierta -- hay una contradiccion sin resolver entre PENDIENTES.md /
% CHECKLIST_CUMPLIMIENTO.md (la marcan abierta) y fichas_antecedentes/
% B3_ByteTrack.md (que da por acordado un alcance limitado a conteo
% agregado). Por eso esta Introduccion se redacta a proposito SIN
% comprometerse con una granularidad de salida especifica: los objetivos
% remiten esa decision a la Seccion de Metodologia y a la auditoria de
% marco_etico_legal.tex (Sec. sec:auditoria), que la resuelven cuando el
% frente tecnico la documente. Metodologia, Resultados y Conclusiones
% siguen pendientes por ese motivo y no se tocan en este cambio.
% =====================================================================
\section{Introducción}
\label{sec:introduccion}

\subsection{Contexto del problema}
\label{sec:contexto}

El Aeropuerto Internacional Jorge Chávez, operado por Lima Airport
Partners (LAP), concentra en determinadas franjas horarias y en puntos
específicos del terminal --- filtros de seguridad, migraciones,
mostradores de \emph{check-in}, salas de embarque --- un volumen de
pasajeros considerablemente mayor que el de otras zonas del recinto en el
mismo momento. Conocer dónde y cuándo se concentran las personas es
información operativa de valor para un operador aeroportuario: permite,
en principio, anticipar la apertura de mostradores adicionales, redirigir
el flujo de pasajeros o dimensionar al personal de atención con mayor
precisión que una revisión manual o periódica.

LAP cuenta ya con infraestructura de videovigilancia de circuito cerrado
en el terminal. Este proyecto explora si esa misma infraestructura puede
aprovecharse, mediante técnicas de visión computacional, para producir
una estimación automatizada de la concentración de personas por zona, sin
que ello suponga identificar, individualizar ni caracterizar
biométricamente a los pasajeros que aparecen en la imagen. Esta última
condición no es una preferencia de diseño sino un límite impuesto por la
Ley 29733 de Protección de Datos Personales y su reglamento vigente
(D.S.\ 016-2024-JUS): un dato biométrico que por sí mismo identifica al
titular constituye dato sensible, y el proyecto se plantea evitar ese
tratamiento desde el diseño en lugar de mitigarlo después. La
Sección~\ref{sec:marco} desarrolla en detalle el marco normativo
aplicable, y la Sección~\ref{sec:antecedentes} revisa la literatura
técnica que hace posible, en principio, estimar una aglomeración sin
recurrir a identificación individual.

Conviene precisar el alcance de esta afirmación desde ya, para no
sobredimensionar el problema que el proyecto atiende. El equipo no cuenta,
al momento de escribir este informe, con evidencia documentada por LAP
--- incidentes registrados de aglomeración, tiempos de espera medidos o
quejas de pasajeros --- que cuantifique el problema en el propio
terminal; esa evidencia se ha solicitado al cliente y su ausencia se deja
constancia expresa en \texttt{PENDIENTES.md}. El proyecto parte, por
tanto, de una necesidad operativa plausible y compartida por la literatura
de conteo de multitudes en espacios de alto tránsito, y no de una cifra
propia de incidentes en el Aeropuerto Jorge Chávez.

\subsection{Objetivos}
\label{sec:objetivos}

\subsubsection*{Objetivo general}

Desarrollar y evaluar un sistema de visión computacional capaz de estimar
la concentración de personas en distintas zonas del Aeropuerto
Internacional Jorge Chávez a partir de imágenes de cámaras fijas, sin
recurrir a reconocimiento facial ni a la extracción de rasgos biométricos
identificantes, y en condiciones de cumplimiento con el marco normativo
peruano de protección de datos personales.

\subsubsection*{Objetivos específicos}

\begin{enumerate}[leftmargin=1.6em,itemsep=0.35em]
  \item Revisar el estado del arte en detección, localización y conteo de
        personas en escenas congestionadas, e identificar qué enfoque
        técnico resulta más adecuado para un entorno aeroportuario y para
        las restricciones de privacidad del proyecto (desarrollado en la
        Sección~\ref{sec:antecedentes}).
  \item Determinar el marco de cumplimiento normativo aplicable al
        tratamiento de imágenes de un espacio de uso público bajo gestión
        concesionada, incluyendo el análisis de impacto y los requisitos
        de privacidad desde el diseño que ese tratamiento exige
        (desarrollado en las Secciones~\ref{sec:marco}
        y~\ref{sec:analisis}).
  \item Diseñar e implementar el componente de detección y estimación,
        con la granularidad de salida --- superficie de densidad, conteo
        agregado por zona o localización individual --- que resulte de
        aplicar el principio de proporcionalidad desarrollado en la
        Sección~\ref{sec:directiva} y la síntesis de
        antecedentes de la Sección~\ref{sec:sintesis-antecedentes}. Esta
        decisión se documentará en la Sección de Metodología cuando el
        frente técnico la cierre; hasta entonces queda registrada como
        abierta en \texttt{PENDIENTES.md} (puntos B-01 y D-11).
  \item Evaluar el desempeño del sistema mediante métricas cuantitativas y
        cualitativas, indicando en cada caso sobre qué conjunto de datos
        fue medido cada resultado --- ninguna cifra obtenida sobre
        conjuntos públicos se presentará como desempeño esperable en el
        terminal sin la validación en sitio que la
        Sección~\ref{sec:ant-p2r} señala como condición previa.
  \item Identificar los requisitos de validación adicionales que quedan
        fuera del alcance material de este Capstone --- en particular el
        acceso a datos del propio terminal y la evaluación de sesgo
        demográfico --- y dejarlos documentados como trabajo pendiente en
        lugar de darlos por resueltos (retomado en la
        Sección~\ref{sec:conclusiones}).
\end{enumerate}

\subsection{Alcance y limitaciones del proyecto}
\label{sec:alcance}

\textbf{Alcance funcional.} El proyecto cubre la estimación de la
concentración de personas por zona a partir de imágenes de cámaras fijas.
Queda expresamente fuera de su alcance cualquier forma de identificación,
reidentificación o caracterización individual de las personas detectadas:
el sistema está pensado para responder \emph{cuántas personas hay, y
dónde}, no \emph{quién es cada una}. La Sección~\ref{sec:ant-p2pnet}
señala, sin embargo, que esta distinción es una cuestión de política de
persistencia y no solo de algoritmo: el mismo modelo que localiza personas
puede, si sus salidas se retienen y se asocian entre cuadros, aproximarse
a un seguimiento de trayectorias. Por eso la granularidad de salida que
finalmente se implemente es, como se indica en el objetivo específico 3,
una decisión que debe quedar documentada por escrito antes de cerrar la
Sección de Metodología.

\textbf{Zonas cubiertas.} Las zonas concretas del terminal sobre las que
se entrena y evalúa el sistema dependen de qué cámaras y qué material
audiovisual habilite LAP para el proyecto, dato que a la fecha de este
informe no ha sido confirmado por el cliente (\texttt{PENDIENTES.md},
punto C-09). Donde este informe menciona zonas a título de ejemplo ---
filtros de seguridad, mostradores de \emph{check-in}, salas de embarque
--- lo hace como ilustración del tipo de espacio al que el sistema podría
aplicarse, no como una cobertura ya acordada con el cliente.

\textbf{Fase del proyecto y limitación de dominio.} Esta primera entrega
corresponde a Capstone I: un desarrollo y validación conceptual del
sistema sobre conjuntos de datos públicos de conteo de multitudes, sin
acceso todavía a material propio del Aeropuerto Jorge Chávez. Esto impone
la limitación más importante del proyecto en esta etapa, y conviene
enunciarla sin atenuarla porque condiciona la lectura de todo lo que
sigue: como documenta en detalle la Sección~\ref{sec:ant-p2r}, los propios
autores de la literatura revisada reconocen una brecha de dominio entre
los conjuntos públicos de entrenamiento y un entorno de despliegue nuevo,
y señalan que el desempeño se degrada de forma marcada cuando no media una
etapa de adaptación de dominio. En consecuencia, ninguna cifra de
desempeño que se reporte en la Sección de Resultados debe leerse como una
predicción de exactitud para el Aeropuerto Jorge Chávez: es, en el mejor
de los casos, una cota de referencia sobre datos públicos, sujeta a
validarse en sitio antes de comunicarse a LAP como una expectativa de
desempeño (Sección~\ref{sec:auditoria}).
